\section{An explicit solution of a 2026 Euler recurrence}
\label{sec:ramanujan-recurrence}

The Ramanujan Challenge for AI \cite{RamanujanChallenge2026} asks, in its
Euler-constant problem, for a proof that the quotient of two distinguished
solutions of the following four-term recurrence tends to \(\gamma\).  We
give a self-contained solution.  The finite sums below belong to the
multiple-Laguerre framework developed by Wolfs--Van Assche
\cite{WolfsVanAssche2025}.  What is proved here is the adjacent
specialization, explicit all-index recurrence certificates, an elementary
positive concentration proof, and the flat matrix field.

Put
\begin{align*}
 A_n&=8n^3+51n^2+105n+68,\\
 B_n&=24n^5+337n^4+1833n^3+4818n^2+6092n+2928,\\
 C_n&=(n+2)(n+3)
 (24n^5+273n^4+1150n^3+2154n^2+1635n+268),\\
 D_n&=(n+1)(n+2)^4(n+3)
 (8n^3+75n^2+231n+232),
\end{align*}
and
\begin{equation}\label{eq:challenge-recurrence}
 A_nu_n=B_nu_{n-1}-C_nu_{n-2}+D_nu_{n-3}
 \qquad(n\ge0).
\end{equation}

\begin{theorem}[Closed solution of the Euler recurrence]
\label{thm:challenge-solution}
For \(n\ge-3\), set \(a=n+3\), \(b=n+4=a+1\), and
\begin{equation}\label{eq:Wab}
 W_{a,b}(k)=
 \frac{(a!b!)^2}{(k!)^3(a-k)!(b-k)!}
 \qquad(0\le k\le a).
\end{equation}
Then the two sequences determined by
\[
 (p_{-3},p_{-2},p_{-1})=(0,7,179),
 \qquad
 (q_{-3},q_{-2},q_{-1})=(1,12,306)
\]
and \eqref{eq:challenge-recurrence} have the exact formulas
\begin{align}
 q_n&=b\sum_{k=0}^{a}W_{a,b}(k),
 \label{eq:qn-closed}\\
 p_n&=b\sum_{k=0}^{a}W_{a,b}(k)
 \bigl(3H_k-H_{a-k}-H_{b-k}\bigr)+a!.
 \label{eq:pn-closed}
\end{align}
In particular,
\begin{equation}\label{eq:challenge-limit}
 \boxed{\displaystyle\lim_{n\to\infty}\frac{p_n}{q_n}=\gamma.}
\end{equation}
\end{theorem}

\begin{proof}
Substitution of \((a,b)=(0,1),(1,2),(2,3)\) gives the six initial
values.  The two symbolic telescoping identities in
\cref{app:telescopers} prove \eqref{eq:challenge-recurrence} for both
closed forms at every \(n\ge0\).  We prove the limit directly.

After cancelling factors independent of \(k\), define a probability law by
\begin{equation}\label{eq:challenge-pmf}
 \Pr(K_n=k)=
 \frac{\binom ak\binom bk/k!}
 {\sum_{j=0}^{a}\binom aj\binom bj/j!}.
\end{equation}
Then
\begin{equation}\label{eq:challenge-expectation}
 \frac{p_n}{q_n}
 =\mathbb E\bigl(3H_{K_n}-H_{a-K_n}-H_{b-K_n}\bigr)
 +\frac{a!}{q_n}.
\end{equation}
The ratio of consecutive unnormalized masses is
\begin{equation}\label{eq:challenge-ratio}
 r_k=\frac{(a-k)(b-k)}{(k+1)^3},
\end{equation}
which decreases strictly in \(k\).  Put \(x_n=(ab)^{1/3}\).  For each
fixed \(\eta>0\), the ratios outside
\([(1-\eta)x_n,(1+\eta)x_n]\) are uniformly separated from one once
\(n\) is large.  Multiplying the ratios on either side of the mode gives
\begin{equation}\label{eq:challenge-concentration}
 \Pr\bigl(|K_n/x_n-1|>\eta\bigr)
 =O\bigl(e^{-c_\eta x_n}\bigr).
\end{equation}
Hence \(K_n/x_n\to1\) in probability.

On the event \(|K_n/x_n-1|\le\eta\), the estimate
\(H_m=\log m+\gamma+O(m^{-1})\) yields
\begin{align*}
 3H_k-H_{a-k}-H_{b-k}
 &=\gamma+\log\frac{k^3}{(a-k)(b-k)}+o(1)\\
 &=\gamma+O(\eta)+o(1).
\end{align*}
The expression is \(O(\log n)\) on the full support, while the exceptional
probability in \eqref{eq:challenge-concentration} is exponentially small in
\(x_n\).  First let \(n\to\infty\), and then let \(\eta\downarrow0\); its
expectation therefore tends to \(\gamma\).  Finally the
\(k=0\) summand in \eqref{eq:qn-closed} gives
\[
 0\le\frac{a!}{q_n}\le\frac1{b\,b!}\longrightarrow0.
\]
Equation \eqref{eq:challenge-expectation} proves
\eqref{eq:challenge-limit}.
\end{proof}

\begin{remark}[What has and has not been solved]
The theorem gives an independent proof of the publicly stated limit problem.
The official proof and code are now public \cite{RamanujanOfficial2026}; the
sentence in Version 2.0 saying otherwise is superseded.  No priority claim is
made for the official Meijer--\(G\) trajectory, Miller selection, or the broad
multiple-Laguerre family.  The comparison with the terminating state is
established in \cref{sec:meijer-duality}.
The convergence \(p_n/q_n\to\gamma\) does not by itself imply the
irrationality of \(\gamma\).
\end{remark}

The same recurrence has a second nontrivial Ap\'ery limit.  Let \(w_n\) be
the solution with
\begin{equation}\label{eq:w-initial}
 (w_{-3},w_{-2},w_{-1})=(1,18,440),
\end{equation}
and let \(\mathfrak h\) be the positive constant in
\eqref{eq:h-constant} below.

\begin{theorem}[Complete initial-value limit classification]
\label{thm:ram-full-solution}
One has
\begin{equation}\label{eq:w-limit}
 \boxed{\displaystyle\frac{w_n}{q_n}\longrightarrow\mathfrak h
 =1.4322057346532244148\ldots.}
\end{equation}
For arbitrary rational initial data
\((x,y,z)=(u_{-3},u_{-2},u_{-1})\), the corresponding solution satisfies
\begin{align}
 \lim_{n\to\infty}\frac{u_n}{q_n}
 =\frac1{136}\bigl[&(-6x+179y-7z)\mathfrak h\notag\\
 &+(-228x-134y+6z)\gamma
 +(142x-179y+7z)\bigr].
 \label{eq:all-initial-limits}
\end{align}
Consequently the set of all rational-initial-data limits is exactly
\begin{equation}\label{eq:attainable-limit-span}
 \boxed{\Q+\Q\gamma+\Q\mathfrak h.}
\end{equation}
No assertion that this span has dimension three is made.
\end{theorem}

\begin{proof}
The official cyclic gauge at zero has rows
\begin{equation}\label{eq:U0-rows}
 \mathsf U(0)=
 \begin{pmatrix}1&18&440\\0&7&179\\0&1&45\end{pmatrix}.
\end{equation}
They generate solutions \(w,p,s\), respectively, and
\(q=w-p+s\).  The official Miller-selection theorem, transported through
the exact gauge of \cref{thm:ram-gauge}, gives
\begin{equation}\label{eq:Miller-limit-vector}
 \frac1{q_n}(w_n,p_n,s_n)
 \longrightarrow
 (\mathfrak h,\gamma,1+\gamma-\mathfrak h);
\end{equation}
see \cite[\S5.2]{RamanujanOfficial2026} and
\cref{thm:meijer-probability}.  The identity \(q=w-p+s\) normalizes the
right side to one.  Solving the initial vector
\((x,y,z)\) in the basis \((q,p,w)\) gives the three coefficients in
\eqref{eq:all-initial-limits}.  Conversely \(q,p,w\) realize
\(1,\gamma,\mathfrak h\), proving \eqref{eq:attainable-limit-span}.
\end{proof}

\begin{theorem}[Complete integer initial lattice]
\label{thm:ram-integer-lattice}
For a solution of \eqref{eq:challenge-recurrence},
\begin{equation}\label{eq:integer-lattice}
 \boxed{\displaystyle
 u_n\in\Z\ \text{for every }n\ge-3
 \Longleftrightarrow
 x,y,z\in\Z\ \text{and }13x+6y+z\equiv0\pmod{17}.}
\end{equation}
\end{theorem}

\begin{proof}
The first new term is
\begin{equation}\label{eq:u0-congruence}
 u_0=\frac{2784x-402y+732z}{17},
\end{equation}
so integrality forces the stated congruence.  Conversely its kernel in
\(\Z^3\) has basis
\[
 (1,0,4),\qquad(0,1,11),\qquad(0,0,17).
\]
The finite companion-gauge certificate in
\cref{app:integer-lattice-certificate} proves that the three corresponding
solutions are integral at every index.  Every integral vector satisfying the
congruence is their integral combination.
\end{proof}

\section{A flat conservative matrix field and exterior arithmetic}
\label{sec:cmf}

For integers \(a,b\ge0\), define
\begin{equation}\label{eq:Fab}
 F_{a,b}(z)={}_{2}F_{2}
 \!\left(\begin{matrix}-a,-b\\1,1\end{matrix};z\right),
 \qquad
 \theta=z\frac d{dz},
 \qquad
 \mathbf Y_{a,b}=
 \begin{pmatrix}F_{a,b}\\ \theta F_{a,b}\\ \theta^2F_{a,b}\end{pmatrix}.
\end{equation}
It satisfies
\begin{equation}\label{eq:Fab-ode}
 \theta^3F_{a,b}=z(\theta-a)(\theta-b)F_{a,b}.
\end{equation}
At \(z=1\), its terminating series gives the exact link
\begin{equation}\label{eq:q-Fab-link}
 F_{a,b}(1)=\sum_{k=0}^{a}\binom ak\binom bk\frac1{k!},
 \qquad q_n=b\,a!b!\,F_{a,b}(1)\quad(a=n+3,\ b=n+4).
\end{equation}

\begin{theorem}[Flat \(SL_3\) conservative matrix field]
\label{thm:flat-cmf}
For \(a,b>0\), put
\begin{align}
 M_a(a,b;z)&=\frac1a
 \begin{pmatrix}
 a&-1&0\\
 0&a&-1\\
 -zab&z(a+b)&a-z
 \end{pmatrix},\label{eq:Ma}\\
 M_b(a,b;z)&=\frac1b
 \begin{pmatrix}
 b&-1&0\\
 0&b&-1\\
 -zab&z(a+b)&b-z
 \end{pmatrix}.\label{eq:Mb}
\end{align}
Then
\[
 \mathbf Y_{a-1,b}=M_a(a,b;z)\mathbf Y_{a,b},
 \qquad
 \mathbf Y_{a,b-1}=M_b(a,b;z)\mathbf Y_{a,b},
\]
and
\begin{align}
 \det M_a&=\det M_b=1,\label{eq:cmf-det}\\
 M_b(a-1,b;z)M_a(a,b;z)
 &=M_a(a,b-1;z)M_b(a,b;z).
 \label{eq:cmf-flatness}
\end{align}
Thus the denominator system of \cref{thm:challenge-solution} is the
adjacent-diagonal path \((a,b)=(n+3,n+4)\) in an explicit flat \(SL_3\)
conservative matrix field.
\end{theorem}

\begin{proof}
The two contiguous relations follow by comparing the terminating
hypergeometric coefficients and using \eqref{eq:Fab-ode} for the third row.
Expansion of the determinants gives \(1\).  Direct matrix multiplication
gives the polynomial identity \eqref{eq:cmf-flatness}; an entrywise symbolic
certificate is also included with the source package.
\end{proof}

\begin{remark}[Path rigidity]
Flatness says that changing a lattice path while keeping the same endpoint
\((a,b)\) leaves \(F_{a,b}\), and hence the endpoint-defined denominator,
unchanged.  The numerator and quotient in \cref{thm:challenge-solution} are
also endpoint-defined by their closed formulas.  Rerouting a path therefore
does not change these particular values.  Changing the endpoint curve is a
different problem and is not excluded by this observation.
\end{remark}

The numerator is the first Frobenius jet of the same flat connection.  Let
\begin{equation}\label{eq:Phi-epsilon}
 \Phi_{a,b}(\epsilon,z)=
 \frac{a!b!}{(1+\epsilon)_a(1+\epsilon)_b}
 {}_3F_3\!\left(
 \begin{matrix}-a-\epsilon,-b-\epsilon,1\\
 1-\epsilon,1-\epsilon,1-\epsilon
 \end{matrix};z\right).
\end{equation}

\begin{theorem}[Flat first-jet \(SL_6\) lift]
\label{thm:ram-sl6}
Fix positive integers \(a,b\), expand the analytic \(\epsilon\)-germ at
zero modulo \(\epsilon^2\), and put
\(\mathbf Y_{a,b}(\epsilon)
=(\Phi_{a,b},\theta\Phi_{a,b},\theta^2\Phi_{a,b})^{\mathsf T}\).
Then
\begin{equation}\label{eq:epsilon-contiguous}
 \mathbf Y_{a-1,b}(\epsilon)=\mathscr M_a(\epsilon)\mathbf Y_{a,b}(\epsilon),
 \qquad
 \mathbf Y_{a,b-1}(\epsilon)=\mathscr M_b(\epsilon)\mathbf Y_{a,b}(\epsilon),
\end{equation}
where
\begin{equation}\label{eq:epsilon-Ma}
 \mathscr M_a(\epsilon)=\frac1a
 \begin{pmatrix}
 a+\epsilon&-1&0\\
 0&a+\epsilon&-1\\
 -zab-\epsilon z(a+b)&z(a+b)+2\epsilon z&a-z-2\epsilon
 \end{pmatrix},
\end{equation}
and
\begin{equation}\label{eq:epsilon-Mb}
 \mathscr M_b(\epsilon)=\frac1b
 \begin{pmatrix}
 b+\epsilon&-1&0\\
 0&b+\epsilon&-1\\
 -zab-\epsilon z(a+b)&z(a+b)+2\epsilon z&b-z-2\epsilon
 \end{pmatrix}.
\end{equation}
In
\(\Q(a,b,z)[\epsilon]/(\epsilon^2)\), these matrices have determinant one
and satisfy the same zero-curvature identity as
\eqref{eq:cmf-flatness}.

Writing
\(\mathscr M_a=\mathscr M_{a,0}+\epsilon\mathscr M_{a,1}\), the block matrix
\begin{equation}\label{eq:SL6-block}
 \begin{pmatrix}
 \mathscr M_{a,0}&0\\
 \mathscr M_{a,1}&\mathscr M_{a,0}
 \end{pmatrix}\in SL_6
\end{equation}
transports
\(\binom{\mathbf Y_{a,b}(0)}
{\partial_\epsilon\mathbf Y_{a,b}(0)}\), and similarly in the
\(b\)-direction.  Thus denominator and numerator path-independence are the
zero- and first-jet parts of one flat connection.

At \(b=a+1,z=1\),
\begin{equation}\label{eq:Phi-pq-jets}
 b\,a!b!\Phi_{a,b}(0,1)=q_{a-3},\qquad
 b\,a!b!\partial_\epsilon\Phi_{a,b}(0,1)=p_{a-3}.
\end{equation}
The boundary \(a=0,b=1\) in \eqref{eq:Phi-pq-jets} is verified directly.
\end{theorem}

\begin{proof}
Modulo \(\epsilon^2\), the Frobenius deformation satisfies
\begin{equation}\label{eq:Phi-dual-ODE}
 (\theta-\epsilon)^3\Phi_{a,b}
 =z(\theta-a-\epsilon)(\theta-b-\epsilon)\Phi_{a,b}.
\end{equation}
Coefficient comparison gives the first two rows of
\eqref{eq:epsilon-contiguous}, while \eqref{eq:Phi-dual-ODE} gives the third.
Direct determinant expansion and multiplication give determinant one and
flatness in the dual-number ring.  Expanding a
dual-number matrix action gives the block form \eqref{eq:SL6-block}.
At \(\epsilon=0\), \(\Phi_{a,b}=F_{a,b}\).  The derivative of the terminating
part gives the harmonic weight in \eqref{eq:pn-closed}; the regularized term
at \(k=b=a+1\) is \(1/(bb!)\) before multiplication by \(b\,a!b!\), and
hence gives exactly the endpoint \(+a!\).
\end{proof}

\begin{proposition}[First-jet exactness barrier]
\label{prop:first-jet-barrier}
Let \(\mathscr K_a\) be the rational carrier defined in
\eqref{eq:Kas-def}.  For \(b=a+1,z=1\),
\begin{equation}\label{eq:Phi-K-contact}
 \Phi_{a,b}(\epsilon,1)-\mathscr K_a(\epsilon)=O(\epsilon^2).
\end{equation}
Return here to the full analytic \(\epsilon\)-germ.  For every \(r\ge1\),
the term with \(k=b+r\) contributes
\begin{equation}\label{eq:E-tail-second-jet}
 [\epsilon^2]\,\Phi_{a,b}(\epsilon,1)\big|_{k=b+r}
 =-\frac{a!b!\,r!(r-1)!}{(b+r)!^3}\ne0.
\end{equation}
The total second-order difference cannot cancel:
\begin{equation}\label{eq:second-jet-total-positive}
 [\epsilon^2]\{\Phi_{a,b}(\epsilon,1)-\mathscr K_a(\epsilon)\}
 =
 \frac1{b^2a!}\left\{
 3H_b-\sum_{r\ge1}\frac{r!(r-1)!}{(b+1)_r^3}\right\}>0.
\end{equation}
Indeed \((b+1)_r\ge(r+1)!\), so the sum is less than one, while
\(3H_b\ge3\).
Thus the first derivative is exactly terminating after its resonant endpoint,
whereas the second derivative develops a genuine infinite hypergeometric
tail.  This is a barrier for this Frobenius deformation, not a universal
no-go for all higher-jet constructions.
\end{proposition}

\begin{proof}
For \(0\le k\le a\), the hypergeometric summand is identically the
corresponding summand of \(\mathscr J_a(\epsilon)\).  At the resonant index
\(k=b\), its analytic continuation is
\[
 \frac{\epsilon}{b^2a!}
 \prod_{j=1}^b\left(1-\frac{\epsilon}{j}\right)^{-3};
\]
its linear coefficient is \(1/(bb!)\), exactly the correction in
\(\mathscr K_a\), and its quadratic coefficient is
\(3H_b/(b^2a!)\).  For \(k=b+r\), \(r\ge1\), direct expansion gives
\eqref{eq:E-tail-second-jet}.  Summing those coefficients gives
\eqref{eq:second-jet-total-positive}; the displayed comparison proves its
strict positivity.  Hence the zero and first jets agree and the second
jets do not.
\end{proof}

\begin{proposition}[Forced divisor and exact Casoratian]
\label{prop:challenge-arithmetic}
Let \(d_b=\lcm(1,\ldots,b)\).  For every \(n\ge-3\),
\begin{equation}\label{eq:challenge-gcd}
 \boxed{\displaystyle
 \frac{b!}{d_b}\gcd\!\left(b,\frac{d_b}{b}\right)
 \mid\gcd(p_n,q_n).}
\end{equation}
If \(u^{(1)},u^{(2)},u^{(3)}\) are three solutions of
\eqref{eq:challenge-recurrence}, put, for \(n\ge-1\),
\[
 \mathcal W_n=\det\bigl(u^{(i)}_{n-j}\bigr)_{1\le i\le3,\,0\le j\le2},
\]
so in particular
\(\mathcal W_{-1}=\det(u^{(i)}_{-1-j})_{1\le i\le3,\,0\le j\le2}\).
Then, for every \(n\ge0\),
\begin{equation}\label{eq:challenge-casoratian}
 \boxed{\displaystyle
 \mathcal W_n=\mathcal W_{-1}
 \frac{A_{n+1}}{136}
 (n+1)!(n+2)!^4(n+3)!.}
\end{equation}
\end{proposition}

\begin{proof}
Write
\[
 W_{a,b}(k)=a!b!\binom ak\binom bk\frac1{k!}.
\]
Put \(g_b=\gcd(b,d_b/b)\).  After division by \(b!/d_b\), each summand of
\(q_n\) is \(b\) times an integer, because \(a!/k!\in\Z\).  The harmonic
part of \(p_n\) is likewise \(b\) times an integer, since
\(d_b(3H_k-H_{a-k}-H_{b-k})\in\Z\).  Its endpoint becomes \(d_b/b\).
Hence every term of both quotients is divisible by \(g_b\), proving
\eqref{eq:challenge-gcd}.

The companion determinant of \eqref{eq:challenge-recurrence} is
\(D_n/A_n\).  Since
\[
 D_n=(n+1)(n+2)^4(n+3)A_{n+1},
\]
iteration from \(-1\) to \(n\) telescopes and gives
\eqref{eq:challenge-casoratian}.
\end{proof}

\begin{remark}[Exterior-volume barrier]
The flat base field has determinant one, whereas scalar elimination together
with its factorial gauge produces the factorial-six Casoratian
\eqref{eq:challenge-casoratian}.  The forced divisor
\eqref{eq:challenge-gcd} is a lower bound for the common divisor, not an
upper bound.  These formulas record the factorial cost but do not rule out
further common-divisor cancellation; no small primitive integer linear form
is obtained here.
\end{remark}

\section{Exact Meijer--\texorpdfstring{\(G\)}{G}/hypergeometric duality}
\label{sec:meijer-duality}

The official solution of the Challenge is now public
\cite{RamanujanOfficial2026}.  It selects a recessive Meijer--\(G\) state by
Miller's algorithm, whereas \cref{sec:ramanujan-recurrence,sec:cmf} uses a
terminating \({}_2F_2\) state and a positive saddle.  The terminating module
is rationally gauge-equivalent to the contragredient of the normalized
official rank-three module.

Write
\[
 F_a(z)=F_{a,a+1}(z),\qquad
 \mathbf Y_a=
 \begin{pmatrix}F_a(1)\\ \theta F_a(1)\\ \theta^2F_a(1)\end{pmatrix}.
\]
The adjacent transition obtained by composing the two contiguous matrices
of \cref{thm:flat-cmf} is
\begin{equation}\label{eq:Delta-adjacent}
 \mathbf Y_a=\mathsf D_a\mathbf Y_{a+1},
\end{equation}
where
\begin{equation}\label{eq:D-adjacent}
\mathsf D_a=
\begin{pmatrix}
1&-\dfrac{2a+3}{(a+1)(a+2)}&\dfrac1{(a+1)(a+2)}\\[2mm]
1&\dfrac{a^2+a-1}{(a+1)(a+2)}&-\dfrac2{a+2}\\[2mm]
-(2a+1)&\dfrac{5a^2+11a+5}{(a+1)(a+2)}&\dfrac{a-2}{a+2}
\end{pmatrix}.
\end{equation}
For \(a\ge1\) this is the composite contiguous transfer.  The boundary
identity \(\mathbf Y_0=\mathsf D_0\mathbf Y_1\) is checked directly from the
terminating series.

Define
\begin{equation}\label{eq:official-Meijer-function}
 \mathcal H_a^G(z)=
 G_{2,3}^{3,2}\!\left(
 z\,\middle|\begin{matrix}1,1\\a+1,a+1,a+1\end{matrix}\right)
\end{equation}
with the Mellin--Barnes normalization of
\cite[\S5.2]{RamanujanOfficial2026}, and let
\[
 \mathbf S_a^G=
 \bigl(\mathcal H_a^G(1),\theta\mathcal H_a^G(1),
       \theta^2\mathcal H_a^G(1)\bigr)
\]
be the official row state, with the normalization of
\cite[\S5.2]{RamanujanOfficial2026}.  With \(b=a+1\), its transfer matrix is
\begin{equation}\label{eq:official-transfer}
 \mathbf S_{a+1}^G=\mathbf S_a^G\mathsf M_a^G,\qquad
 \mathsf M_a^G=
 \begin{pmatrix}
 0&b^3&b^3(3a+4)\\
 0&-3b^2&-b^2(8a+11)\\
 1&3b&6a^2+12a+5
 \end{pmatrix}.
\end{equation}
Put
\begin{equation}\label{eq:official-normalized-gauge}
 \mathsf P_a^G=b^{-2}\mathsf M_a^G,\qquad
 \mathsf B_a=
 \begin{pmatrix}
 b^2&-(2a+1)&1\\
 0&a/b&-1/b\\
 0&1/b&0
 \end{pmatrix}.
\end{equation}

\begin{theorem}[Exact rational \(SL_3\) coboundary]
\label{thm:ram-gauge}
For every \(a\ge0\),
\begin{equation}\label{eq:ram-determinants}
 \det\mathsf D_a=\det\mathsf B_a=\det\mathsf P_a^G=1,
 \qquad \det\mathsf M_a^G=(a+1)^6,
\end{equation}
and the exact gauge identity is
\begin{equation}\label{eq:ram-coboundary}
 \boxed{\displaystyle
 \mathsf D_a^{-\mathsf T}\mathsf B_{a+1}
 =\mathsf B_a(\mathsf P_a^G)^{-\mathsf T}.}
\end{equation}

If \(\mathbf V_a=\mathbf Y_a^{\mathsf T}\mathsf B_a\) and
\(\Pi_a^G=\mathsf M_0^G\cdots\mathsf M_{a-1}^G\), then
\begin{align}
 \mathbf V_{a+1}&=\mathbf V_a(\mathsf P_a^G)^{-\mathsf T},
 \label{eq:dual-trajectory}\\
 \boxed{\displaystyle
 \mathbf Y_a^{\mathsf T}\mathsf B_a
 =(a!)^2(1,-1,1)(\Pi_a^G)^{-\mathsf T}.}
 \label{eq:dual-trajectory-product}
\end{align}
The denominator coordinate has the precise factorial gauge
\begin{equation}\label{eq:q-dual-coordinate}
 \boxed{\displaystyle
 q_{a-3}=(a!)^2\mathbf V_ae_1.}
\end{equation}
Equivalently, \(\widehat{\mathbf V}_a=(a!)^2\mathbf V_a\) has first
coordinate \(q_{a-3}\) and evolves by the denominator-cleared dual matrix
\begin{equation}\label{eq:denominator-cleared-dual}
 \widehat{\mathbf V}_{a+1}
 =\widehat{\mathbf V}_a\widehat{\mathsf M}_a^G,\qquad
 \widehat{\mathsf M}_a^G=b^4(\mathsf M_a^G)^{-\mathsf T},\qquad
 \det\widehat{\mathsf M}_a^G=b^6.
\end{equation}
\end{theorem}

\begin{proof}
For \(a\ge1\), equation \eqref{eq:Delta-adjacent} is the product
\(M_b(a,a+2;1)M_a(a+1,a+2;1)\) from
\cref{thm:flat-cmf}; multiplication gives \eqref{eq:D-adjacent}.  At
\(a=0\), direct evaluation of the two terminating polynomials gives the
same identity.
The determinants in \eqref{eq:ram-determinants} are direct expansions.
Substitution of \cref{eq:D-adjacent,eq:official-normalized-gauge} into
\eqref{eq:ram-coboundary} reduces its nine entries to rational identities in
\(a\).

Rearranging \eqref{eq:ram-coboundary} gives
\(\mathsf D_a^{\mathsf T}\mathsf B_a
=\mathsf B_{a+1}(\mathsf P_a^G)^{\mathsf T}\), which proves
\eqref{eq:dual-trajectory}.  Since
\(\mathbf Y_0=(1,0,0)^{\mathsf T}\) and
\(\mathbf Y_0^{\mathsf T}\mathsf B_0=(1,-1,1)\), iteration gives
\eqref{eq:dual-trajectory-product}.  Finally,
\(\mathbf V_ae_1=b^2F_a(1)\), whereas
\eqref{eq:q-Fab-link} gives
\(q_{a-3}=(a!)^2b^2F_a(1)\).  This proves
\eqref{eq:q-dual-coordinate}; the last assertions follow by rescaling.
\end{proof}

The gauge fixes the finite connection pairing, not only its projective
direction.

\begin{theorem}[Finite Meijer--\(G\)/\({}_2F_2\) bilinear identity]
\label{thm:ram-bilinear}
For every \(a\ge0\),
\begin{equation}\label{eq:ram-bilinear-matrix}
 \boxed{\displaystyle
 \mathbf S_a^G\mathsf B_a^{\mathsf T}\mathbf Y_a=(a!)^2.}
\end{equation}
Writing \(b=a+1\), \(F=F_a(1)\), and
\(H=\mathcal H_a^G(1)\), this is
\begin{align}
 b^2HF
 &+(\theta H)\left(-(2a+1)F+\frac ab\theta F
                         +\frac1b\theta^2F\right)\notag\\
 &+(\theta^2H)\left(F-\frac1b\theta F\right)=(a!)^2.
 \label{eq:ram-bilinear-expanded}
\end{align}
Thus the normalized conserved primal--dual pairing is exactly one.
\end{theorem}

\begin{proof}
The official normalization gives
\(\mathbf S_0^G(1,-1,1)^{\mathsf T}=1\).  Use
\(\mathbf S_a^G=\mathbf S_0^G\Pi_a^G\) and
\eqref{eq:dual-trajectory-product}; the factors \(\Pi_a^G\) cancel.
Expanding \(\mathsf B_a^{\mathsf T}\mathbf Y_a\) gives
\eqref{eq:ram-bilinear-expanded}.
\end{proof}

The determinant-six phenomenon in \cref{prop:challenge-arithmetic} can now
be located exactly.  If \(\mathsf U(a)\) is the cyclic gauge of the official
solution, then
\begin{equation}\label{eq:cyclic-gauge-det}
 \det\mathsf U(a)=(a+1)^5(a+2)A_a,\qquad \det\mathsf U(0)=136.
\end{equation}
The coboundary transformation
\begin{equation}\label{eq:companion-gauge-definition}
 \mathsf C^{\mathrm{comp}}(a)=
 \mathsf U(a)^{-1}\widehat{\mathsf M}_a^G\mathsf U(a+1)
\end{equation}
gives the scalar companion transfer.  Consequently
\begin{align}
 \det\bigl(\mathsf C^{\mathrm{comp}}(0)\cdots
                 \mathsf C^{\mathrm{comp}}(n)\bigr)
 &=((n+1)!)^6\frac{\det\mathsf U(n+1)}{136}\notag\\
 &=\frac{A_{n+1}}{136}(n+1)!(n+2)!^4(n+3)!,
 \label{eq:casoratian-gauge-origin}
\end{align}
which is the multiplier of \(\mathcal W_{-1}\) in
\eqref{eq:challenge-casoratian}.  The sixth factorial power comes from the
denominator-cleared dual determinant; the shifted factorials and
\(A_{n+1}/136\) are the endpoint cyclic-gauge determinant.

\section{Euler--Beta--log transmutation and the offset ladder}
\label{sec:beta-log}

The adjacent sums are an exact integral transform of the cubic Pilehrood
polynomials, not merely an analogous family.  Define
\begin{align}
 \mathcal Q_m(t)&=\sum_{j=0}^m\binom mj^3j!\,t^j,
 \label{eq:cubic-Q}\\
 \mathcal P_m(t)&=\sum_{j=0}^m\binom mj^3j!
   \bigl(3H_{m-j}-2H_j\bigr)t^j.
 \label{eq:cubic-P}
\end{align}
These are the cubic polynomials occurring in the Euler-constant construction
of Hessami Pilehrood--Hessami Pilehrood
\cite{Pilehrood2013,PilehroodBell2012}.  For \(d\ge1\), put
\begin{equation}\label{eq:Beta-operator}
 \mathcal B_d[f]=\frac1{(d-1)!}\int_0^1(1-t)^{d-1}f(t)\,dt
\end{equation}
and
\begin{equation}\label{eq:transmutation-operator}
 \mathcal T_{m,d}[f](z)=
 \frac{(m+d)!}{m!^2}z^m\mathcal B_d[f(t/z)].
\end{equation}

\begin{theorem}[Exact Beta--log transmutation]
\label{thm:beta-log-transform}
For every \(m\ge0\) and \(d\ge1\), the first identity below holds for
\(z\ne0\) and extends polynomially to \(z=0\); at \(z=1\), the second
identity also holds:
\begin{align}
 F_{m,m+d}(z)&=\mathcal T_{m,d}[\mathcal Q_m](z),
 \label{eq:offset-F-transform}\\
 \mathcal G_{m,d}(1)&=
 \mathcal T_{m,d}[\mathcal P_m+\mathcal Q_m\log t](1),
 \label{eq:offset-G-transform}
\end{align}
where
\begin{equation}\label{eq:Gmd-def}
 \mathcal G_{m,d}(z)=
 \sum_{k=0}^m\binom mk\binom{m+d}k\frac{z^k}{k!}
 \bigl(3H_k-H_{m-k}-H_{m+d-k}\bigr).
\end{equation}
Moreover the transform intertwines Euler derivations:
\begin{equation}\label{eq:Beta-intertwining}
 \theta_z\mathcal T_{m,d}
 =\mathcal T_{m,d}(m-\theta_t).
\end{equation}

For the Challenge specialization \(m=a\), \(d=1\), \(z=1\),
\begin{align}
 q_{a-3}&=(a+1)^3a!\int_0^1\mathcal Q_a(t)\,dt,
 \label{eq:q-Beta-transform}\\
 p_{a-3}&=(a+1)^3a!
 \int_0^1\bigl(\mathcal P_a(t)+\mathcal Q_a(t)\log t\bigr)\,dt+a!.
 \label{eq:p-Beta-transform}
\end{align}
\end{theorem}

\begin{proof}
Put \(j=m-k\).  The beta integrals
\begin{align*}
 \mathcal B_d[t^j]&=\frac{j!}{(j+d)!},\\
 \mathcal B_d[t^j\log t]&=\frac{j!}{(j+d)!}(H_j-H_{j+d})
\end{align*}
turn the coefficient of \(z^{m-j}=z^k\) on the right of
\eqref{eq:offset-F-transform} into
\[
 \frac{m!(m+d)!}{k!^3(m-k)!(m+d-k)!},
\]
which is the coefficient of \(F_{m,m+d}\).  Adding
\(3H_{m-j}-2H_j\) and \(H_j-H_{j+d}\) gives exactly
\(3H_k-H_{m-k}-H_{m+d-k}\), proving
\eqref{eq:offset-G-transform} at \(z=1\).  Differentiating
\(z^mf(t/z)\) gives
\eqref{eq:Beta-intertwining}.  Finally
\cref{eq:q-Beta-transform,eq:p-Beta-transform} follow from
\cref{eq:qn-closed,eq:pn-closed}.
\end{proof}

\begin{corollary}[Fixed-offset Euler ladder]
\label{cor:offset-ladder}
For every fixed \(d\ge1\),
\begin{equation}\label{eq:offset-gamma-limit}
 \boxed{\displaystyle
 \frac{\mathcal G_{m,d}(1)}{F_{m,m+d}(1)}\longrightarrow\gamma.}
\end{equation}
\end{corollary}

\begin{proof}
Use the probability weights
\[
 \Prob(K=k)\propto\binom mk\binom{m+d}k\frac1{k!}.
\]
Their consecutive ratio is
\((m-k)(m+d-k)/(k+1)^3\), so, with
\(s_m=(m(m+d))^{1/3}\), the same ratio argument as in
\cref{thm:challenge-solution} gives \(K/s_m\to1\) exponentially in
probability.  On the window \(|K/s_m-1|\le\eta\),
\[
 3H_K-H_{m-K}-H_{m+d-K}
 =\gamma+\log\frac{K^3}{(m-K)(m+d-K)}+o(1)
 =\gamma+O(\eta)+o(1).
\]
The logarithmic growth on the complement is absorbed by the exponential
tail.  First let \(m\to\infty\), and then let \(\eta\downarrow0\).
Taking expectations proves \eqref{eq:offset-gamma-limit}.
\end{proof}

The multiple-Laguerre background and general approximation families are
already present in the cited literature and in
\cite{WolfsVanAssche2025}.  The assertion here is the exact transmutation
between that cubic system and the 2026 Challenge normalization, together
with its full fixed-offset form.

\section{Rational Gamma jets and contact polynomials}
\label{sec:gamma-jets}

The numerator and denominator are the first two jets of one rational
deformation.  Continue to write \(b=a+1\), and define
\begin{align}
 \mathscr J_a(s)&=a!b!\sum_{k=0}^a
 \frac1{(1-s)_k^3(1+s)_{a-k}(1+s)_{b-k}},
 \label{eq:Jas-def}\\
 \mathscr K_a(s)&=\mathscr J_a(s)+\frac{s}{b\,b!}.
 \label{eq:Kas-def}
\end{align}
The small final term is essential: it records the resonant endpoint
correction \(+a!\) in \eqref{eq:pn-closed}.

\begin{theorem}[Master Gamma-jet transform]
\label{thm:gamma-jet}
For every \(a\ge0\),
\begin{equation}\label{eq:zero-first-jets}
 \boxed{\displaystyle
 q_{a-3}=b\,a!b!\,\mathscr K_a(0),\qquad
 p_{a-3}=b\,a!b!\,\mathscr K_a'(0),}
\end{equation}
and hence
\begin{equation}\label{eq:gamma-first-log-jet}
 \frac{p_{a-3}}{q_{a-3}}=(\log\mathscr K_a)'(0).
\end{equation}
As \(a\to\infty\), locally uniformly on the strip \(|\Re s|<1\),
\begin{equation}\label{eq:Gamma-master-limit}
 \boxed{\displaystyle
 \frac{\mathscr J_a(s)}{\mathscr J_a(0)},\quad
 \frac{\mathscr K_a(s)}{\mathscr K_a(0)}
 \longrightarrow
 \Gamma(1-s)^3\Gamma(1+s)^2.}
\end{equation}
Consequently, for every \(m\ge2\),
\begin{equation}\label{eq:zeta-log-jets}
 \boxed{\displaystyle
 \frac{\left.\dfrac{d^m}{ds^m}\log\mathscr K_a(s)\right|_{s=0}}
 {(3+2(-1)^m)(m-1)!}\longrightarrow\zeta(m),}
\end{equation}
while the first logarithmic jet tends to \(\gamma\).
Every finite-stage logarithmic jet in \eqref{eq:zeta-log-jets} is rational.
\end{theorem}

\begin{proof}
Differentiating a reciprocal Pochhammer symbol at zero gives
\begin{align*}
 \left.\frac d{ds}\log(1-s)_k^{-3}\right|_{s=0}&=3H_k,\\
 \left.\frac d{ds}\log(1+s)_r^{-1}\right|_{s=0}&=-H_r.
\end{align*}
Equations \cref{eq:qn-closed,eq:pn-closed} now give
\eqref{eq:zero-first-jets}; the derivative of \(s/(bb!)\) produces exactly
the endpoint \(a!\).

Normalize the summands of \(\mathscr J_a(0)\) to the probability law
\eqref{eq:challenge-pmf}.  Uniformly for \(s\) in a compact subset of the
strip \(|\Re s|<1\), the ratio of a deformed summand to its value at zero is
 \begin{align}
 &\left(\frac{k!}{(1-s)_k}\right)^3
 \frac{(a-k)!}{(1+s)_{a-k}}
 \frac{(b-k)!}{(1+s)_{b-k}}\notag\\
 &\quad=\Gamma(1-s)^3\Gamma(1+s)^2
 k^{3s}(a-k)^{-s}(b-k)^{-s}(1+o(1)).
 \label{eq:uniform-gamma-ratio}
 \end{align}
Choose a shrinking central window
\(|k/(ab)^{1/3}-1|\le\eta_a\), where
\(\eta_a\downarrow0\) and \(\eta_a^2(ab)^{1/3}\to\infty\).
The concentration estimate \eqref{eq:challenge-concentration} makes its
complement exponentially small.  On the window,
\(k=(ab)^{1/3}(1+o(1))\), \(a-k=a(1+o(1))\), and
\(b-k=b(1+o(1))\); hence the power factor in
\eqref{eq:uniform-gamma-ratio} tends uniformly to one.  Uniform Gamma-ratio
bounds on each compact substrip give at most a fixed polynomial in
\(k+1,a-k+1,b-k+1\), which the exponential tail absorbs.  This proves the
locally uniform limit for \(\mathscr J_a\).  The extra term in
\(\mathscr K_a\) is \(O((bb!)^{-1})\), whereas
\(\mathscr J_a(0)\ge1\), so the same limit holds for \(\mathscr K_a\).

Finally,
\begin{equation}\label{eq:Gamma-log-series}
 \log\{\Gamma(1-s)^3\Gamma(1+s)^2\}
 =\gamma s+\sum_{m\ge2}\frac{3+2(-1)^m}{m}\zeta(m)s^m.
\end{equation}
Local uniformity permits termwise differentiation of logarithms near zero
and proves \eqref{eq:zeta-log-jets}.  Rationality at finite \(a\) follows
because \(\mathscr K_a\in\Q(s)\) and \(\mathscr K_a(0)\ne0\).
\end{proof}

For example, as \(a\to\infty\), the correctly normalized first cases are
\begin{equation}\label{eq:first-zeta-jets}
 \frac{(\log\mathscr K_a)''(0)}5\to\zeta(2),\qquad
 \frac{(\log\mathscr K_a)^{(3)}(0)}2\to\zeta(3),\qquad
 \frac{(\log\mathscr K_a)^{(4)}(0)}{30}\to\zeta(4).
\end{equation}
The raw derivative \(\mathscr K_a''(0)\) cannot replace the second
logarithmic derivative.

\begin{corollary}[Exact parity separation]
\label{cor:gamma-jet-parity}
As \(a\to\infty\), locally uniformly on \(|\Re s|<1\),
\begin{align}
 \frac{\mathscr K_a(s)\mathscr K_a(-s)}{\mathscr K_a(0)^2}
 &\longrightarrow
 \left(\frac{\pi s}{\sin\pi s}\right)^5,
 \label{eq:even-jet-limit}\\
 \frac{\mathscr K_a(s)}{\mathscr K_a(-s)}
 &\longrightarrow\frac{\Gamma(1-s)}{\Gamma(1+s)}.
 \label{eq:odd-jet-limit}
\end{align}
Thus the logarithm of the first quotient generates only even zeta values,
while the logarithm of the second generates \(\gamma\) and the odd zeta
values.
\end{corollary}

\begin{proof}
Multiply and divide the two limits in \eqref{eq:Gamma-master-limit}, and use
\(\Gamma(1-s)\Gamma(1+s)=\pi s/\sin\pi s\).
\end{proof}

The Bell-polynomial zeta tower itself is classical
\cite{PilehroodBell2012}.  The point of
\cref{thm:gamma-jet,cor:gamma-jet-parity} is its exact finite connection to
the individual 2026 Challenge approximants, including their endpoint term.

The same deformation has a finite integer spectral avatar.  Put
\begin{equation}\label{eq:contact-denominator}
 \mathscr D_a(s)=
 \prod_{j=1}^a(s-j)^3(s+j)^2(s+a+1).
\end{equation}

\begin{theorem}[Monic integer contact polynomial]
\label{thm:contact-polynomial}
There is a unique polynomial \(\mathscr N_a\in\Z[s]\) such that
\begin{equation}\label{eq:contact-rational-form}
 \mathscr K_a(s)=
 \frac{\mathscr N_a(s)}{a!(a+1)^2\mathscr D_a(s)}.
\end{equation}
It is monic of degree \(5a+2\), and
\begin{equation}\label{eq:contact-degree}
 \deg\{\mathscr N_a-s\mathscr D_a\}=3a.
\end{equation}
Consequently \(\mathscr N_a\) and \(s\mathscr D_a\) have the same first
\(2a+1\) Newton power sums.  Explicitly, if the zeros of
\(\mathscr N_a\) are counted with multiplicity, then for
\(1\le m\le2a+1\),
\begin{equation}\label{eq:contact-Newton}
 \sum_{\mathscr N_a(\nu)=0}\nu^m
 =3\sum_{j=1}^aj^m+2\sum_{j=1}^a(-j)^m+(-a-1)^m.
\end{equation}
At the opposite end of the polynomial,
\begin{align}
 \mathscr N_a(0)
 &=(-1)^a(a+1)(a!)^4q_{a-3},
 \label{eq:contact-N0}\\
 \mathscr N_a'(0)
 &=(-1)^a(a+1)(a!)^4
 \left[p_{a-3}-q_{a-3}\left(H_a-\frac1{a+1}\right)\right].
 \label{eq:contact-N1}
\end{align}
\end{theorem}

\begin{proof}
Multiplication of \eqref{eq:Kas-def} by
\(a!(a+1)^2\mathscr D_a(s)\) clears every Pochhammer denominator.  Each
resulting quotient is in \(\Z[s]\); the endpoint term becomes exactly
\(s\mathscr D_a(s)\).  The latter is monic of degree \(5a+2\), while the
\(k\)-th summand of the remaining part has degree \(3a-k\).  The \(k=0\)
term has nonzero leading coefficient, proving monicity and
\eqref{eq:contact-degree}.

Two monic degree-\(5a+2\) polynomials whose difference has degree \(3a\)
share their first \(2a+1\) coefficients below the leading term.  Newton's
identities give \eqref{eq:contact-Newton}; the roots of
\(s\mathscr D_a\) are \(0\), three copies of \(1,\ldots,a\), two copies of
\(-1,\ldots,-a\), and \(-a-1\).  Finally evaluate
\eqref{eq:contact-rational-form} and its logarithmic derivative at zero,
use \eqref{eq:zero-first-jets}, and compute
\[
 \frac{\mathscr D_a'(0)}{\mathscr D_a(0)}
 =-H_a+\frac1{a+1}.
\]
This gives \cref{eq:contact-N0,eq:contact-N1}.
\end{proof}

\section{Probabilistic realizations and zero geometry}
\label{sec:ram-probability}

The positive law in \eqref{eq:challenge-pmf} has two complementary
realizations: exponential maxima at finite index and independent Bernoulli
variables through the zeros of its generating polynomial.

For \(m\ge1\), let \(M_m\) be the maximum of \(m\) independent
\(\operatorname{Exp}(1)\) variables, put \(M_0=0\), and use independent
copies whenever superscripts are present.  Then
\begin{equation}\label{eq:maxima-mgf}
 \mathbb Ee^{tM_m}=\prod_{j=1}^m\frac j{j-t}
 =\frac{m!}{(1-t)_m}\qquad(\Re t<1).
\end{equation}

\begin{theorem}[Exact finite maxima model]
\label{thm:ram-probability}
For \(a\ge0\), let \(J_a\) have the law
\begin{equation}\label{eq:Ja-law}
 \Prob(J_a=k)=
 \frac{\binom ak\binom{a+1}k/k!}{F_{a,a+1}(1)}
 \qquad(0\le k\le a),
\end{equation}
and, conditionally on \(J_a\), define
\begin{equation}\label{eq:Za0}
 Z_a^{(0)}=
 M_{J_a}^{(1)}+M_{J_a}^{(2)}+M_{J_a}^{(3)}
 -M_{a-J_a}^{(4)}-M_{a+1-J_a}^{(5)}.
\end{equation}
Then
\begin{equation}\label{eq:exact-maxima-mgf}
 \boxed{\displaystyle
 \mathbb Ee^{tZ_a^{(0)}}
 =\frac{\mathscr J_a(t)}{\mathscr J_a(0)}}
 \qquad(|\Re t|<1).
\end{equation}
If
\begin{equation}\label{eq:Za-correction}
 c_a=\frac{a!}{q_{a-3}},\qquad Z_a=Z_a^{(0)}+c_a,
\end{equation}
then
\begin{equation}\label{eq:finite-mean-approximant}
 \boxed{\displaystyle\mathbb EZ_a=\frac{p_{a-3}}{q_{a-3}}.}
\end{equation}
Moreover, as \(a\to\infty\), with independent standard Gumbel variables
\(G_1,\ldots,G_5\),
\begin{equation}\label{eq:Za-all-moments}
 Z_a\longrightarrow X_0:=G_1+G_2+G_3-G_4-G_5
\end{equation}
in distribution and with convergence of every moment.  Equivalently,
locally uniformly for \(|\Re t|<1\),
\begin{equation}\label{eq:Za-mgf-limit}
 \mathbb Ee^{tZ_a}\longrightarrow
 \Gamma(1-t)^3\Gamma(1+t)^2.
\end{equation}
Thus
\begin{align}
 \kappa_1(Z_a)&\longrightarrow\gamma,\notag\\
 \kappa_m(Z_a)&\longrightarrow
 (m-1)!\{3+2(-1)^m\}\zeta(m)\qquad(m\ge2).
 \label{eq:Za-cumulants}
\end{align}
\end{theorem}

\begin{proof}
Multiply the weight in \eqref{eq:Ja-law} by the five conditional moment
generating functions from \eqref{eq:maxima-mgf}.  All ordinary factorials
cancel, leaving exactly the summand of \(\mathscr J_a(t)\); normalization at
zero proves \eqref{eq:exact-maxima-mgf}.  Taking the first derivative and
using \eqref{eq:zero-first-jets} gives
\eqref{eq:finite-mean-approximant}.  Since \(c_a\to0\),
\cref{thm:gamma-jet} gives \eqref{eq:Za-mgf-limit} on a neighborhood of the
origin.  Uniform exponential moments on every smaller strip imply
distributional and all-moment convergence.  Finally differentiate the
logarithm and use \eqref{eq:Gamma-log-series}.
\end{proof}

The distinction between the two deformations is important:
\begin{equation}\label{eq:mgf-not-K}
 \mathbb Ee^{tZ_a}
 =e^{c_at}\frac{\mathscr J_a(t)}{\mathscr J_a(0)}
 \ne\frac{\mathscr K_a(t)}{\mathscr K_a(0)}
\end{equation}
at finite \(a\).  The function \(\mathscr K_a\) is the exact rational
zero/first-jet carrier; \(\mathscr J_a\) is the exact probability carrier.
They have the same Gamma limit.

We next identify the full official Meijer state.  Let
\begin{equation}\label{eq:Xn-def}
 A_{n,i}\sim\Gamma(n+1,1)\quad(1\le i\le3),\qquad
 E_4,E_5\sim\operatorname{Exp}(1)
\end{equation}
be independent, and put
\begin{equation}\label{eq:Xn-ratio}
 X_n=\log\frac{E_4E_5}{A_{n,1}A_{n,2}A_{n,3}}.
\end{equation}

\begin{theorem}[Stop-loss interpretation of the Meijer state]
\label{thm:meijer-probability}
For the official normalization in \cref{sec:meijer-duality},
\begin{equation}\label{eq:Meijer-stop-loss}
 \boxed{\displaystyle
 \mathcal H_n^G(e^{-x})=(n!)^3\mathbb E(X_n-x)_+}
 \qquad(x\in\R).
\end{equation}
Consequently
\begin{equation}\label{eq:Meijer-state-probability}
 \boxed{\displaystyle
 \frac{\bigl(\mathcal H_n^G(1),\theta\mathcal H_n^G(1),
 \theta^2\mathcal H_n^G(1)\bigr)}{(n!)^3}
 =\bigl(\mathbb E(X_n)_+,\Prob(X_n>0),f_{X_n}(0)\bigr).}
\end{equation}
The functional shift is
\begin{equation}\label{eq:Meijer-functional-shift}
 \mathcal H_{n+1}^G(z)=z\theta^2\mathcal H_n^G(z),
\end{equation}
and hence
\begin{equation}\label{eq:probability-conservation}
 f_{X_n}(x)=e^x(n+1)^3\mathbb E(X_{n+1}-x)_+.
\end{equation}

At \(n=0\),
\begin{equation}\label{eq:h-constant}
 \mathfrak h:=\mathbb E(X_0)_+
 =\int_0^\infty e^{-x}\frac{\log x}{x-1}\,dx
 =1.4322057346532244148\ldots,
\end{equation}
whereas
\begin{equation}\label{eq:Powers-gamma}
 \mathbb EX_0=\Prob(X_0>0)=\gamma.
\end{equation}
In particular \(\mathfrak h\ne\gamma\).
\end{theorem}

\begin{proof}
The Mellin transform of \(X_n\) is
\begin{equation}\label{eq:Xn-mgf}
 \mathbb Ee^{sX_n}
 =\frac{\Gamma(1+s)^2\Gamma(n+1-s)^3}{(n!)^3},
 \qquad -1<\Re s<n+1.
\end{equation}
Perron's inverse Mellin formula on a line \(0<c<n+1\) gives
\eqref{eq:Meijer-stop-loss}; differentiation in \(x\), using
\(\theta=-d/dx\) at \(z=e^{-x}\), gives
\eqref{eq:Meijer-state-probability}.  Shifting the three gamma factors in
the Mellin--Barnes integral proves \eqref{eq:Meijer-functional-shift}, and
comparison with \eqref{eq:Meijer-stop-loss} gives
\eqref{eq:probability-conservation}.  At \(n=0\), substitute the five
exponential densities into \(\mathbb E(X_0)_+\), integrate four variables
first, and apply Frullani's identity; the remaining integral is
\[
 \int_0^\infty e^{-x}\frac{\log x}{x-1}\,dx,
\]
which proves \eqref{eq:h-constant}.
The probability identity in \eqref{eq:Powers-gamma} is due to Powers
\cite{Powers2023}; the expectation identity is immediate from
\(\mathbb EG_j=\gamma\).
\end{proof}

The Bernoulli realization comes from real-rootedness.  For integers
\(b\ge a\ge1\), write the zeros of \(F_{a,b}\), with multiplicities, as
\(-\rho_{a,b,1},\ldots,-\rho_{a,b,a}\).

\begin{theorem}[Poisson--binomial law and zero-density limit]
\label{thm:poisson-binomial}
All zeros of \(F_{a,b}\) are negative real numbers, and
\begin{equation}\label{eq:Poisson-binomial-factor}
 \frac{F_{a,b}(z)}{F_{a,b}(1)}
 =\prod_{j=1}^a(1-p_j+p_jz),
 \qquad p_j=\frac1{1+\rho_{a,b,j}}.
\end{equation}
Thus the law with weights proportional to
\(\binom ak\binom bk/k!\) is exactly a sum of independent Bernoulli
variables.  Its mass sequence is strictly log-concave and ultra-log-concave,
and it is monotone in each of \(a,b\) in the monotone-likelihood-ratio order.

Fix \(d\ge0\), put \(b=a+d\) and \(s_a=(ab)^{1/3}\).  As \(a\to\infty\),
if \(J_a\) has these weights, then
\begin{align}
 \mathbb EJ_a&\sim s_a,&
 \Var(J_a)&\sim\frac{s_a}{3},
 \label{eq:Ja-mean-var}\\
 \frac{J_a-\mathbb EJ_a}{\sqrt{\Var(J_a)}}&\Longrightarrow N(0,1).
 \label{eq:Ja-CLT}
\end{align}
The corresponding local central limit theorem holds, and \(J_a/s_a\)
satisfies a large-deviation principle with speed \(s_a\) and rate
\begin{equation}\label{eq:Ja-rate}
 I(x)=
 \begin{cases}
 3(x\log x-x+1),&x\ge0,\\
 +\infty,&x<0,
 \end{cases}
 \qquad(0\log0:=0).
\end{equation}

Finally, the locally finite zero measures converge vaguely on
\([0,\infty)\):
\begin{equation}\label{eq:zero-density-limit}
 \boxed{\displaystyle
 \frac1{s_a}\sum_{j=1}^a\delta_{\rho_{a,b,j}}
 \Longrightarrow
 \frac{\sqrt3}{2\pi}\rho^{-2/3}\,d\rho.}
\end{equation}
\end{theorem}

\begin{proof}
The polynomial \(L_a(-z)=\sum_k\binom akz^k/k!\) has only negative real
zeros.  The coefficient sequence \(\{\binom bk\}_{k\ge0}\) is a classical
P\'olya--Schur multiplier sequence because its exponential generating
function is \(L_b(-z)\); see \cite{PolyaSchur1914}.  Applying the multiplier
theorem to \(L_a(-z)\)
proves real-rootedness, and positive coefficients make every root negative.
Factorization and normalization at \(z=1\) give
\eqref{eq:Poisson-binomial-factor}.

The consecutive mass ratio
\[
 \frac{w_{k+1}}{w_k}=\frac{(a-k)(b-k)}{(k+1)^3}
\]
is strictly decreasing, proving strict log-concavity.  Ultra-log-concavity
follows because \(w_k/\binom ak=\binom bk/k!\) has decreasing consecutive
ratio.  Finally
\[
 \frac{w_{a+1,b;k}}{w_{a,b;k}}=\frac{a+1}{a+1-k}
\]
is increasing in \(k\), and similarly for \(b\), proving the likelihood
ratio assertion.

A Taylor expansion about the unique real saddle
\(\kappa_a\), determined by
\[
 (a-\kappa_a)(b-\kappa_a)=(\kappa_a+1)^3,
\]
shows that the discrete law has one mode or two adjacent modes neighboring
\(\kappa_a\).  Its logarithmic curvature is
\[
 \frac1{a-\kappa_a}+\frac1{b-\kappa_a}
 +\frac3{\kappa_a+1}\sim\frac3{s_a},
\]
which gives \eqref{eq:Ja-mean-var}.  The Bernoulli representation,
\(\Var(J_a)\to\infty\), and the standard Poisson--binomial local limit
theorem \cite[Chapter~VII]{Petrov1975} give the central and local limit
assertions.  Stirling's
formula at \(k=xs_a\), uniformly on compact \(x\)-intervals, gives
\eqref{eq:Ja-rate} and exponential tightness gives the full large-deviation
principle.

For \(z>0\), tilt the weights by \(z^k\).  Their ratio has the extra factor
\(z\), so \(J_{a,z}/s_a\to z^{1/3}\).  Therefore
\begin{equation}\label{eq:Stieltjes-transform-limit}
 \frac1{s_a}\frac{F'_{a,b}(z)}{F_{a,b}(z)}
 =\frac{\mathbb E_zJ_{a,z}}{s_az}\longrightarrow z^{-2/3}.
\end{equation}
The left side is the Stieltjes transform of the measure on the left of
\eqref{eq:zero-density-limit}, whereas
\[
 \int_0^\infty\frac1{\rho+z}
 \frac{\sqrt3}{2\pi}\rho^{-2/3}\,d\rho=z^{-2/3}.
\]
Uniqueness and the standard continuity theorem for Stieltjes transforms
\cite[Chapter~I]{Akhiezer1965} give the vague limit.
\end{proof}

No simplicity or strict-interlacing assertion is needed above.  Those
stronger statements require a separate strict multiplier or
multiple-orthogonality input and are not claimed here.

\begin{remark}[Priority boundary for the Ramanujan package]
The official solution already proves Question~2 through a rank-three
Meijer--\(G\) trajectory, a Mellin--Barnes recessive estimate, a rational
companion gauge, and inverse-transpose Miller selection
\cite{RamanujanOfficial2026}.  The cubic sums and their Bell--zeta
approximants occur in \cite{PilehroodBell2012,Pilehrood2013}, while the
five-exponential ratio law, its signed-Gumbel logarithm, and the associated
Gamma-product transform occur in \cite{Powers2023}.  The general
multiple-orthogonal and multiple-Laguerre background is developed in
\cite{VanAsscheWolfs2023,WolfsVanAssche2025}.  The contribution of the
present sections is the explicit finite-level transmutation among the
Challenge recurrence, its terminating \({}_2F_2\) dual, the official
Meijer--\(G\) trajectory, and the cubic Bell family, together with the exact
finite probability carrier and the all-order logarithmic-jet limits in the
Challenge normalization.
\end{remark}

\section{The regularized Gamma determinant}
\label{sec:gamma-determinant}

The positive trace-class operator of \cref{sec:spectral} encodes the single
Euler trace \(S_2\).  A different, signed Hilbert--Schmidt operator encodes
the complete limiting Gamma jet.  On
\(\ell^2(\N)^{\oplus5}\), define
\begin{equation}\label{eq:A-Gamma}
 \mathsf A_\Gamma=\diag\left(
 \frac1n,\frac1n,\frac1n,-\frac1n,-\frac1n:n\ge1\right).
\end{equation}
Let \(\mathsf A_{\Gamma,N}\) be its compression to the first \(N\)
coordinates in each of the five summands.  We use the standard
Hilbert--Schmidt determinant
\[
 \det{}_2(I-T)=\det\bigl((I-T)e^T\bigr).
\]

\begin{theorem}[Regularized determinant and zeta traces]
\label{thm:gamma-det2}
One has
\(\mathsf A_\Gamma\in\mathcal S_2\setminus\mathcal S_1\), and
\begin{equation}\label{eq:FP-trace-gamma}
 \boxed{\displaystyle
 \FPTr\mathsf A_\Gamma
 :=\lim_{N\to\infty}
 \left(\Tr\mathsf A_{\Gamma,N}-\log N\right)=\gamma.}
\end{equation}
For every \(m\ge2\),
\begin{equation}\label{eq:Gamma-operator-traces}
 \boxed{\displaystyle
 \Tr\mathsf A_\Gamma^m
 =\{3+2(-1)^m\}\zeta(m).}
\end{equation}
Equivalently,
\begin{equation}\label{eq:Gamma-det2}
 \boxed{\displaystyle
 \Gamma(1-t)^3\Gamma(1+t)^2
 =e^{\gamma t}\det{}_2(I-t\mathsf A_\Gamma)^{-1}.}
\end{equation}
In particular
\begin{equation}\label{eq:Gamma-even-odd-traces}
 \Tr\mathsf A_\Gamma^{2r}=5\zeta(2r),\qquad
 \Tr\mathsf A_\Gamma^{2r+1}=\zeta(2r+1)\qquad(r\ge1).
\end{equation}
\end{theorem}

\begin{proof}
Square summability and failure of absolute summability are immediate from
the spectrum.  The truncated trace is \(H_N\), so
\eqref{eq:FP-trace-gamma} is Euler's definition.  Taking powers of the five
eigenvalue strings gives \eqref{eq:Gamma-operator-traces}.  Finally, for a
Hilbert--Schmidt operator,
\[
 \log\det{}_2(I-t\mathsf A_\Gamma)
 =-\sum_{m\ge2}\frac{t^m}{m}\Tr\mathsf A_\Gamma^m.
\]
Compare with \eqref{eq:Gamma-log-series}; this first proves
\eqref{eq:Gamma-det2} for \(|t|<1\), and meromorphic continuation gives the
identity on \(\C\).
\end{proof}

This regularized determinant is an explicit diagonal spectral realization of the
Gamma product, but it does not furnish a new irrationality result for odd
zeta values.  It must also not be confused with the ordinary Fredholm
determinant \(\Delta_2\) of the positive trace-class operator in
\cref{thm:Delta2}.

\section{Cross-place Euler--Pad\'e and continued-fraction convergence}
\label{sec:cross-place-pade}

The Pad\'e system for Euler's factorial series gives a separate arithmetic
transfer.  Put
\begin{equation}\label{eq:Euler-Pade-Q}
 Q_n(T)=\sum_{h=0}^nh!\binom nh^2(-T)^h,
\end{equation}
and let \(P_n(T)\), of degree at most \(n-1\), be the Pad\'e numerator
characterized by
\begin{equation}\label{eq:Euler-Pade-condition}
 Q_n(T)\sum_{m\ge0}m!T^m-P_n(T)=O(T^{2n})
\end{equation}
as a formal series.  For a prime \(p\), the series
\begin{equation}\label{eq:p-adic-Euler-series}
 E_p(t)=\sum_{m\ge0}m!t^m
\end{equation}
converges for \(|t|_p\le1\).  At a real negative argument put
\begin{equation}\label{eq:real-Euler-function}
 E_\infty(t)=\int_0^\infty\frac{e^{-x}}{1-tx}\,dx.
\end{equation}
The standard Pad\'e remainder gives, for \(t\in\Z_p\),
\begin{equation}\label{eq:p-adic-remainder-bound}
 v_p\!\left(E_p(t)-\frac{P_n(t)}{Q_n(t)}\right)
 \ge2v_p(n!)-v_p(Q_n(t)),
\end{equation}
whenever \(Q_n(t)\ne0\); see
\cite{ErnvallMatalaSeppala2022}.

\begin{theorem}[Unit-boundary residue criterion]
\label{thm:p-adic-unit-criterion}
Let \(p\) be prime and \(t\in\Z_p^\times\).  For
\(0\le r<p\), define
\begin{equation}\label{eq:q-rp}
 q_{r,p}(X)=\sum_{h=0}^rh!\binom rh^2(-X)^h\in\mathbb F_p[X].
\end{equation}
If
\begin{equation}\label{eq:good-residue-condition}
 q_{r,p}(\bar t)\ne0\qquad(0\le r<p),
\end{equation}
then every \(Q_n(t)\) is a \(p\)-adic unit and
\begin{equation}\label{eq:full-p-adic-convergence}
 \boxed{\displaystyle
 \frac{P_n(t)}{Q_n(t)}\longrightarrow E_p(t).}
\end{equation}
The following are sample good nonzero residue classes:
\begin{equation}\label{eq:good-residue-table}
 \begin{array}{c|c}
 p&t\pmod p\\ \hline
 3&2\\
 5&2,3,4\\
 7&4,5\\
 11&2,3,5,7\\
 13&6,9,11,12
 \end{array}.
\end{equation}
In particular, at \(t=-1\), the full convergent sequence tends to
\(E_p(-1)\) for \(p=2,3,5,13\).
\end{theorem}

\begin{proof}
Modulo \(p\), all terms with \(h\ge p\) vanish and, for \(h<p\),
\(\binom nh\bmod p\) depends only on \(n\bmod p\).  Hence
\begin{equation}\label{eq:Q-period-p}
 Q_n(t)\equiv q_{r,p}(\bar t)\pmod p
 \qquad(n\equiv r\pmod p).
\end{equation}
Condition \eqref{eq:good-residue-condition} makes every denominator a unit;
then \eqref{eq:p-adic-remainder-bound} tends to infinity.  The table is a
finite calculation in the indicated residue fields.

For \(p=2,t=-1\), an elementary calculation modulo four gives
\begin{equation}\label{eq:Q-v2}
 v_2(Q_n(-1))=
 \begin{cases}0,&n\text{ even},\\1,&n\text{ odd}.
 \end{cases}
\end{equation}
The same remainder estimate therefore proves full convergence also at two.
For \(p=3,5,13\), the assertion is contained in the table.
\end{proof}

The congruence behind a simultaneous subsequence is stronger than periodicity.

\begin{theorem}[A common real and all-prime subsequence]
\label{thm:common-placewise-subsequence}
For all integers \(n\ge0\) and \(k\ge1\),
\begin{equation}\label{eq:Carlitz-congruence}
 Q_{n+k}(T)\equiv Q_n(T)\pmod{k\Z[T]}.
\end{equation}
Let \(N_j=\lcm(1,\ldots,j)\).  For every fixed negative integer \(t\), the
single rational sequence
\begin{equation}\label{eq:common-rational-sequence}
 x_j(t)=\frac{P_{N_j}(t)}{Q_{N_j}(t)}
\end{equation}
satisfies
\begin{equation}\label{eq:common-placewise-limits}
 \boxed{\displaystyle
 x_j(t)\longrightarrow E_\infty(t)\quad\text{in }\R,
 \qquad
 x_j(t)\longrightarrow E_p(t)\quad\text{in }\Q_p
 \text{ for every fixed }p.}
\end{equation}
At \(t=-1\), the real limit is the Euler--Gompertz constant
\(\delta=eE_1(1)\), while the \(p\)-adic limit is \(E_p(-1)\).
\end{theorem}

\begin{proof}
For each \(h\), the polynomial
\(h!\binom Xh=X(X-1)\cdots(X-h+1)\) belongs to \(\Z[X]\).  Hence
\[
 h!\left\{\binom{n+k}h^2-\binom nh^2\right\}
 =\left\{h!\binom{n+k}h-h!\binom nh\right\}
  \left\{\binom{n+k}h+\binom nh\right\}
\]
is divisible by \(k\), coefficient by coefficient in
\eqref{eq:Euler-Pade-Q}.  This proves \eqref{eq:Carlitz-congruence}.
Taking \(n=0,k=N_j\) gives
\begin{equation}\label{eq:QN-unit-all-p}
 Q_{N_j}(T)\equiv1\pmod{N_j\Z[T]}.
\end{equation}
Thus for every fixed \(p\) and all \(j\ge p\), \(Q_{N_j}(t)\) is a
\(p\)-adic unit, and \eqref{eq:p-adic-remainder-bound} proves the \(p\)-adic
limit.  The real Stieltjes--Pad\'e convergence to
\eqref{eq:real-Euler-function} holds for every \(t<0\), and therefore along
the same subsequence.
\end{proof}

\begin{proposition}[\(p\)-adic denominator interpolation]
\label{prop:p-adic-denominator-interpolation}
For every prime \(p\), the series
\begin{equation}\label{eq:q-p-interpolation}
 q_p(x)=\sum_{h\ge0}h!\binom xh^2\qquad(x\in\Z_p)
\end{equation}
converges uniformly and defines a \(1\)-Lipschitz function satisfying
\(q_p(n)=Q_n(-1)\) for \(n\in\N_0\).  Moreover
\begin{equation}\label{eq:q-p-root-criterion}
 \boxed{\displaystyle
 \sup_{n\ge0}v_p(Q_n(-1))=\infty
 \Longleftrightarrow q_p\text{ has a zero in }\Z_p.}
\end{equation}
In particular, if \(q_p\) is rootless then the full Euler convergent
sequence at \(t=-1\) tends to \(E_p(-1)\).
\end{proposition}

\begin{proof}
Uniform convergence follows from \(v_p(h!)\to\infty\) and
\(\binom xh\in\Z_p\).  For each summand,
\[
 h!\left\{\binom xh^2-\binom yh^2\right\}
 =\left(x^{\underline h}-y^{\underline h}\right)
  \left\{\binom xh+\binom yh\right\},
\]
which is divisible by \(x-y\); passing to the uniform limit gives the
Lipschitz bound.  The interpolation identity follows because
\(\binom nh=0\) for \(h>n\).  If the valuations are unbounded, compactness
of \(\Z_p\) supplies a convergent subsequence of the corresponding integers,
and continuity gives a zero.  Conversely, approximate a zero by
nonnegative integers and use continuity.  If there is no zero, compactness
bounds \(v_p(q_p(x))\), and \eqref{eq:p-adic-remainder-bound} proves the
last assertion.
\end{proof}

The word ``simultaneous'' in \cref{thm:common-placewise-subsequence} is
placewise: no convergence in the restricted adelic product topology is
asserted.  Nor does the theorem prove irrationality of any \(E_p(-1)\).
At the residue \(n=p-1\),
\begin{equation}\label{eq:Kurepa-contact}
 Q_{p-1}(-1)\equiv\sum_{h=0}^{p-1}h!\pmod p,
\end{equation}
so Kurepa's left-factorial problem appears as one boundary congruence, but is
not solved here.  Full convergence at \(t=-1\) for every prime remains open;
the observed growth bound \(v_p(Q_n(-1))=O(\log n)\) is not presently a
theorem.  The existence of a zero of \(q_p\) means only that denominator
valuations are unbounded; it is not equivalent to failure of convergence,
which would require comparison with \(v_p(n!)\).
